% C.24, strengthened in place; coefficients are genuine Q-vector spaces.
\begin{theorem}[Proved: rational-vector coefficient cocycle colimits]
\label{C4:acyclic-cocycle-colimits}\label{C4:vector-cocycle-colimits}
Let $I$ be a nonempty small filtered category. For each $i\in I$,
let $V_i$ be a countable rational vector space and let $r_i\geq2$.
Allow arbitrary functors in a diagram
$i\mapsto\mathcal A_{r_i}(V_i)$, without requiring that they come
from linear maps of the coefficient spaces. For its categorical
colimit $A$, one has
\[
 \widetilde H_*(BA;\mathbb Q)=0
 \quad\Longleftrightarrow\quad
 A\text{ is exact over }\mathbb Z.
\]
Thus every total, or merely weakly contractible-base, colimit of
this form is integrally exact, including when the coefficient vector
spaces have unbounded or infinite dimension.
\end{theorem}
\begin{proof}
We give the finite obstruction and its functorial control, which
make the argument work in arbitrary coefficient dimension.
For a finite family $f=(a_0,\ldots,a_k):P\to Q$ in
$\mathcal A_r(V)$, let $K_f^*$ be the subcomplex
\[
 K_f^n=\ker\left(
 C^n(\Delta^Q;\mathbb Q)\longrightarrow
 \bigoplus_{i=1}^k C^n(\Delta^P;\mathbb Q)\right),
 \qquad b\longmapsto(a_i^*b-a_0^*b)_{i=1}^k.
 \tag{VC.1}
\]
The prescribed cocycle $c_Q$ belongs to $K_f^r\otimes V$.
Its class
\[
 o_f=[c_Q]\in H^r(K_f^*)\otimes_{\mathbb Q}V
 \tag{VC.2}
\]
vanishes precisely when there is a primitive $b_Q$ with all
$a_i^*b_Q$ equal. Put $b_P=a_0^*b_Q$ in that case. The rational
primitive criterion of Proposition~\ref{C4:rational-primitive-criterion}
works verbatim for $V$: rational affine barycenters of $V$-valued
primitives are defined, and the expansion by arbitrary
$V$-valued $(r-1)$-cochains is the same ordered free-amalgamation
Ramsey class. Hence
\[
 o_f=0
 \quad\Longleftrightarrow\quad
 f\text{ has an integral affine equalizer}.
 \tag{VC.3}
\]
The vector space $H^r(K_f^*)$ is finite dimensional, since $P,Q$
are finite. Consequently (VC.2) involves only finitely many
vectors of $V$, even if $V$ is infinite dimensional.

We record explicitly the cochain comparison, including its
naturality for arbitrary functors. For
$F:\mathcal A_r(V)\to\mathcal A_s(W)$ put
$K_q^F(B)=C_q(\Delta^{\operatorname{vert}F(B)};\mathbb Q)$ and
form the first-quadrant double complex
\[
 \mathcal B_{p,q}(F)=
 \bigoplus_{B_0\to\cdots\to B_p}K_q^F(B_0).
 \tag{VC.4}
\]
Its vertical augmentation is a quasi-isomorphism to the
constant-coefficient bar complex of $\mathcal A_r(V)$, because
each full simplex is nonempty and contractible. The global
target cocycle defines a cochain in bidegree $(0,s)$ of
$\operatorname{Hom}_{\mathbb Q}(\mathcal B(F),W)$; its horizontal
coboundary is zero by naturality. There is a map to the
corresponding double complex for the identity functor of the
target: send each bar string to its $F$-image and retain its
simplex. On vertical augmentations this map is the ordinary
bar map of $F$. In the target identity double complex, every
tuple action is a disjoint union of representables, and the
horizontal degree-zero quotient is the semisimplicial complex
of $K(W,s)$. Its tautological cocycle is the universal
$W$-valued class: on the difference of top simplices with labels
$w$ and $0$ it evaluates to $w$, and their boundaries agree.
Consequently the class of the global cocycle in (VC.4) is
exactly the pullback under $BF$ of this universal class.

If the smooth vertex set $S$ of $F$ has no fixed vertex, every
nonempty increasing tuple in $S$ has a proper open stabilizer.
Theorem~\ref{B:saturation} and
Proposition~\ref{B:cocycle-audit} therefore make every component
of its category of elements contractible. Thus the positive
horizontal homology of (VC.4) vanishes. Its degree-zero
horizontal quotient is $C_*(\Delta^S;\mathbb Q)_G$: every finite
tuple occurs at a finite source stage, and homogeneity realizes
the identifications. Projection to that quotient is a
quasi-isomorphism. Dualizing into $W$ identifies its invariant
cochain cohomology with $H^*(B\mathcal A_r(V);W)$, and the
preceding bar map identifies the cocycle class naturally with
the pulled-back universal class. All total-complex diagonals
are finite, and duality into a rational vector space is exact,
so these comparisons require no convergence or finite-dimension
assumption on $V,W$. If $S$ has a fixed vertex, invariant coning
instead gives a primitive already in bidegree $(0,s-1)$ of
the double cochain complex, so the pullback class is zero.

Consider now any functor
$F:\mathcal A_r(V)\to\mathcal A_r(W)$, and write
\[
 L=(BF)_*:V=H_r(B\mathcal A_r(V);\mathbb Q)
       \longrightarrow H_r(B\mathcal A_r(W);\mathbb Q)=W.
 \tag{VC.5}
\]
We claim that
\[
 F(f)\text{ has an integral affine equalizer}
 \quad\Longleftrightarrow\quad
 (1\otimes L)o_f=0.
 \tag{VC.6}
\]
Use the smooth ordered vertex set $S$ of $F$, with its invariant
$W$-valued cocycle $c_S$, as in Theorem~\ref{C4:degree-colimit-exclusion}.
If $S$ has a fixed vertex, invariant coning makes $c_S$ exact,
and every image family equalizes. The natural bar/simplex
comparison also gives $L=0$, so (VC.6) holds in this case.

Suppose that $S$ has no fixed vertex. The finite-support argument
from Theorem~\ref{C4:degree-colimit-exclusion} applies without
change to $\mathcal A_r(V)$. It supplies an equivariant point
section $j:M\to S$ and the augmentation-preserving support map
\[
 p:\mathbb Q[S]\longrightarrow\mathbb Q[M],
 \qquad p(v)=|B_v|^{-1}\sum_{x\in B_v}x,
 \qquad pj=\operatorname{id}.
\]
For every finite source object $Q$, the support of every vertex
of $F(Q)$ is contained in $Q$. Thus $p$ restricts to a chain map
from the full simplex on $F(Q)$ to that on $Q$, and these maps
are natural under every source embedding. The section $j$ also
restricts naturally: choose one vertex of $F$ of the one-point
object and transport it along its unique embeddings.
It need not preserve the displayed linear orders; the induced
map of alternating simplex chains is nevertheless defined and
natural.

The ordered prism gives an equivariant homotopy $H$ from
$(jp)_*$ to the identity on the simplex chain complex of $S$.
The invariant cochain $j^*c_S$ represents the homomorphism $L$
in
\[
 H^r\bigl(C^*(\Delta^M;W)^G\bigr)
   =H^r(B\mathcal A_r(V);W)
   =\operatorname{Hom}_{\mathbb Q}(V,W).
 \tag{VC.7}
\]
Here the first identification is obtained by the bar/simplex
comparison above for the identity source vertex functor. The second is the
universal coefficient theorem over $\mathbb Q$, since the
classifying space is $K(V,r)$. Naturality follows from the double-complex maps and the
objectwise natural section $j$, whose vertical augmentation
induces the identity on the source bar complex. In particular, for the tautological
$V$-valued source cocycle $c_M$, there is an invariant cochain
$\alpha$ such that
\[
 j^*c_S=Lc_M+d\alpha.
\]
The prism identity therefore gives the global equality
\[
 c_S=p^*(Lc_M)+d\bigl(c_SH+p^*\alpha\bigr).
 \tag{VC.8}
\]
The error primitive in (VC.8) is invariant, so its restrictions
to $F(P),F(Q)$ commute with every arrow in the given family.
If $(1\otimes L)o_f=0$, choose compatible primitives of $Lc_Q$
and pull them back by $p$, adding the error primitive from
(VC.8). This supplies compatible primitives for $F(f)$.
Conversely, pull any compatible primitive for $F(f)$ back by
$j$ and subtract $\alpha$. This supplies a compatible primitive
of $Lc_Q$. This proves (VC.6) by (VC.3).

Two further consequences of the same construction will be used. A functor
$\mathcal A_r(V)\to\mathcal A_s(W)$ with $s\geq|Q|$ sends a
family with target $Q$ to an integrally equalizable family,
by the support-prism rank bound. If $s<r$, the invariant target
cocycle is exact because
$H^s(K(V,r);W)=0$; hence every image family equalizes.
If $V=0$, the source category already has the Ramsey property,
so every image family has an integral affine equalizer. If $W=0$,
the target has that property. Neither case can occur along a
persistent nonequalizable family.

Assume that $BA$ is rationally acyclic, but that $A$ is not
integrally exact. A nonequalizable finite family is represented
at some stage $i$ and remains nonequalizable at every later
stage. Its fixed finite target bounds all future degrees by the
rank bound. Choose a stage $j$ attaining the largest degree in
this filtered tail. No later degree can be smaller, by the
preceding vanishing assertion; hence all degrees on the
cofinal category $j/I$ equal some $r$.

On this tail, filtered categorical colimits commute with nerves
and rational homology, so
\[
 0=H_r(BA;\mathbb Q)=\operatorname*{colim}_{k\in j/I} V_k.
\]
Write the obstruction $o_{f_j}$ in a basis of the finite
dimensional space $H^r(K_{f_j}^*)$. Each of its finitely many
coordinates in $V_j$ becomes zero at some stage. Filteredness
provides one later stage where they all vanish. Apply (VC.6) to
the transition functor from $j$ to that stage: the image family
has an integral affine equalizer, a contradiction.
The converse follows from integral exactness applied to the
constant rational coefficient diagram.
\end{proof}
